(1) 容器绝热且向内部真空膨胀，外部体积功 \(W_{\rm ext}=0\)、热输入 \(Q=0\)，因此整体第一定律 \(\Delta u=0\)。但相互作用势能项 \(-a/v\) 依赖密度，故 \(\Delta u=0\) 不等于 \(\Delta T=0\)。

(2) 代入内能状态函数得 \(C_VT_f-a/v_f=C_VT_i-a/v_i\)，即 \(\boxed{T_f=T_i+(a/C_V)(1/v_f-1/v_i)}\)。因 \(v_f>v_i\)，括号为负，气体冷却；分子吸引的负势能在分离时增加，消耗了部分热运动能量。

(3) 固定 \(u\) 微分 \(0=C_VdT+a v^{-2}dv\)，得 \(\boxed{\mu_J=(\partial T/\partial v)_u=-a/(C_Vv^2)<0}\)。当 \(a\to0\) 时降温消失，恢复经典理想气体自由膨胀的等温结论。此Joule系数与固定焓的Joule–Thomson系数不同。

(4) 由热力学 \(du=Tds-pdv\)，并用 \(du=C_VdT+a dv/v^2\) 与 \(p=RT/(v-b)-a/v^2\)，得到 \(ds=C_VdT/T+Rdv/(v-b)\)。故 \(\boxed{\Delta s=C_V\ln(T_f/T_i)+R\ln[(v_f-b)/(v_i-b)]}\)。在物理稳定单相、\(v_i>b\) 且 \(T_f>0\) 的实际可达自由膨胀中总熵非负；若求得非正温或越过相变，单相模型不能直接适用。
